\chapter{Implementacja aplikacji benchmarkowej}
\label{chap:implementation}

\section{Cel rozdziału}

Celem niniejszego rozdziału jest przedstawienie implementacji aplikacji benchmarkowej wykorzystanej w badaniach porównawczych. Aplikację przygotowano w dwóch wariantach architektonicznych: jako system monolityczny oraz jako system oparty na mikroserwisach. Oba warianty realizują ten sam proces biznesowy i operują na tych samych danych wejściowych oraz wyjściowych. Dzięki temu możliwe było przeprowadzenie rzetelnego porównania wpływu architektury na wydajność, skalowalność oraz koszty utrzymania.

W celu zwiększenia wartości porównawczej eksperymentu każdą z wersji architektonicznych zaimplementowano w trzech stosach technologicznych:
\begin{itemize}
    \item Java z wykorzystaniem frameworka Spring,
    \item JavaScript z wykorzystaniem środowiska Node.js,
    \item Python z wykorzystaniem frameworka Flask.
\end{itemize}

W rezultacie przygotowano sześć implementacji aplikacji: trzy monolity oraz trzy warianty mikroserwisowe. Pozwala to analizować nie tylko wpływ samego stylu architektonicznego, ale również sprawdzić, czy obserwowane zależności są spójne pomiędzy różnymi ekosystemami programistycznymi.

\section{Charakterystyka aplikacji benchmarkowej}

Aplikacja benchmarkowa została osadzona w domenie \textit{image processing}. Jej zadaniem jest wyszukiwanie obrazów podobnych do obrazu referencyjnego na podstawie ekstrakcji cech oraz porównania z wcześniej przygotowanym indeksem wektorów cech. Wynik działania systemu stanowi lista obrazów spełniających zadany próg podobieństwa wraz z odpowiadającymi im metadanymi.

Wybrany scenariusz jest odpowiedni do badań architektonicznych z kilku powodów. Po pierwsze, obejmuje kilka kolejnych etapów przetwarzania, które mogą zostać naturalnie wydzielone do osobnych komponentów. Po drugie, zawiera zarówno element obliczeniowy, związany z ekstrakcją cech z obrazu, jak i etap wyszukiwania podobieństwa oraz pobierania metadanych. Po trzecie, logika ta może zostać zrealizowana zarówno w jednej aplikacji monolitycznej, jak i w formie rozproszonego łańcucha mikroserwisów, bez zmiany sensu biznesowego całego systemu.

W obu wariantach architektonicznych aplikacja korzysta z tego samego zbioru danych wejściowych, obejmującego:
\begin{itemize}
    \item buforowane obrazy referencyjne,
    \item indeks cech obrazów zapisany w pliku \texttt{features.json},
    \item metadane obrazów zapisane w pliku \texttt{metadata.json}.
\end{itemize}

Dane te są ładowane przy starcie aplikacji do struktur pamięciowych, co pozwala ograniczyć wpływ zewnętrznej warstwy persystencji na wyniki eksperymentu. Dzięki temu badanie koncentruje się na wpływie architektury systemu i modelu komunikacji, a nie na kosztach dostępu do bazy danych.

\section{Model wejścia i wyjścia systemu}

Z perspektywy klienta testowego aplikacja udostępnia interfejs HTTP. W przedstawionym scenariuszu badawczym punkt końcowy przyjmuje parametry sterujące procesem wyszukiwania obrazów podobnych. Najważniejszymi parametrami są próg podobieństwa oraz rozmiar jądra wykorzystywanego w etapie ekstrakcji cech.

Przykładowe wywołanie punktu końcowego można przedstawić następująco:

\begin{lstlisting}[language={},caption={Przykładowe wywołanie punktu końcowego HTTP},label={lst:image-search-http}]
GET /images/search?threshold=0.2&kernel=3
\end{lstlisting}

Logiczny model żądania można również zapisać w postaci uproszczonego obiektu JSON:

\begin{lstlisting}[caption={Logiczny model żądania wyszukiwania},label={lst:image-search-request}]
{
  "threshold": 0.2,
  "kernel": 3
}
\end{lstlisting}

W odpowiedzi system zwraca listę obiektów reprezentujących obrazy uznane za podobne do obrazu referencyjnego. Każdy element odpowiedzi zawiera identyfikator obrazu, jego nazwę oraz listę tagów opisujących zawartość. Przykładową odpowiedź przedstawiono w listingu~\ref{lst:image-search-response}.

\begin{lstlisting}[caption={Przykładowa odpowiedź systemu},label={lst:image-search-response}]
[
  {
    "id": "img_001",
    "name": "forest-path",
    "tags": ["forest", "nature", "path"]
  },
  {
    "id": "img_014",
    "name": "green-woods",
    "tags": ["trees", "forest", "landscape"]
  }
]
\end{lstlisting}

Z punktu widzenia eksperymentu istotne jest to, że wszystkie implementacje zwracają ten sam typ odpowiedzi końcowej. Pozwala to uznać badane warianty za funkcjonalnie równoważne.

\section{Logika przetwarzania obrazu}

Proces biznesowy realizowany przez aplikację składa się z kilku kolejnych etapów. Najpierw wybierany jest obraz referencyjny, który w wariancie testowym pobierany jest z wcześniej załadowanego zbioru obrazów. Następnie z obrazu wyznaczany jest wektor cech. W przedstawionej implementacji cechy są obliczane na podstawie uproszczonego przetwarzania obrazu, obejmującego:
\begin{itemize}
    \item konwersję bufora obrazu do tablicy wartości pikseli,
    \item normalizację do skali szarości,
    \item zastosowanie uśredniającego jądra konwolucyjnego o zadanym rozmiarze,
    \item budowę histogramu o ustalonej liczbie przedziałów,
    \item normalizację histogramu do postaci wektora cech.
\end{itemize}

Otrzymany wektor porównywany jest następnie z indeksem cech wszystkich obrazów zapisanych w systemie. Miara podobieństwa została oparta na odległości euklidesowej pomiędzy wektorem zapytania a wektorami zapisanymi w indeksie. Do wyniku wyszukiwania trafiają identyfikatory tych obrazów, dla których odległość nie przekracza zadanego progu podobieństwa.

W ostatnim etapie system pobiera metadane dla znalezionych identyfikatorów obrazów i buduje odpowiedź zwracaną do klienta. Taki przepływ danych jest wspólny dla wariantu monolitycznego i mikroserwisowego. Różnica dotyczy wyłącznie sposobu rozmieszczenia odpowiedzialności pomiędzy komponentami oraz modelu komunikacji wewnętrznej.

\section{Implementacja wariantu monolitycznego}

Wariant monolityczny został zaimplementowany jako pojedyncza aplikacja uruchamiana w jednym procesie. W tej wersji wszystkie etapy przetwarzania obrazu, tj. obsługa żądania HTTP, ekstrakcja cech, wyszukiwanie podobieństwa oraz pobieranie metadanych, realizowane są lokalnie wewnątrz jednego artefaktu wdrożeniowego.

W implementacji Java wariant monolityczny został oparty na frameworku Spring. Aplikacja udostępnia punkt końcowy \texttt{/images/search}, a podczas uruchamiania ładuje obrazy, indeks cech oraz metadane do struktur pamięciowych. Przepływ wykonania obejmuje kolejno: odczyt parametrów wejściowych, wybór obrazu referencyjnego, ekstrakcję cech, porównanie z indeksem oraz mapowanie wyników do odpowiedzi zwracanej klientowi. Analogiczny model działania zastosowano również w implementacjach JavaScript i Python.

Takie podejście upraszcza strukturę systemu. Komunikacja pomiędzy poszczególnymi częściami logiki odbywa się przez lokalne wywołania funkcji lub metod, bez konieczności użycia protokołów sieciowych. Oznacza to niższy narzut komunikacyjny, prostszy model wdrożenia oraz mniejszą liczbę komponentów operacyjnych. Z drugiej strony cały system stanowi jeden artefakt, który w przypadku skalowania poziomego musi być replikowany jako całość.

Schemat infrastruktury wariantu monolitycznego przedstawiono na rysunku~\ref{fig:infra-monolith}, natomiast sekwencję realizacji żądania w tym wariancie pokazano na rysunku~\ref{fig:seq-monolith}.

\begin{figure}[htbp]
    \centering
    \includegraphics[width=0.9\textwidth]{images/architecture/infra-monolit.pdf}
    \caption{Schemat infrastruktury wariantu monolitycznego}
    \label{fig:infra-monolith}
\end{figure}

\begin{figure}[htbp]
    \centering
    \includegraphics[width=0.95\textwidth]{images/architecture/monolit-sequence-diagram.pdf}
    \caption{Diagram sekwencji przetwarzania w wariancie monolitycznym}
    \label{fig:seq-monolith}
\end{figure}

\section{Implementacja wariantu mikroserwisowego}

Wariant mikroserwisowy został zbudowany jako system złożony z pięciu współpracujących usług. Każda z nich odpowiada za wyodrębniony etap przetwarzania, a komunikacja pomiędzy usługami realizowana jest za pomocą gRPC. Zastosowanie tego protokołu miało na celu ograniczenie narzutu komunikacyjnego względem klasycznego modelu HTTP/JSON oraz zapewnienie jednolitego kontraktu między komponentami systemu.

Architektura mikroserwisowa obejmuje następujące komponenty:
\begin{itemize}
    \item \textbf{Gateway} -- warstwa wejściowa udostępniająca interfejs HTTP dla klienta zewnętrznego,
    \item \textbf{Search} -- usługa orkiestrująca cały proces wyszukiwania,
    \item \textbf{Feature} -- usługa odpowiedzialna za ekstrakcję wektora cech z obrazu referencyjnego,
    \item \textbf{Similarity} -- usługa porównująca wektor zapytania z indeksem cech i zwracająca listę identyfikatorów obrazów spełniających próg podobieństwa,
    \item \textbf{Metadata} -- usługa pobierająca metadane obrazów dla identyfikatorów zwróconych przez usługę \textit{Similarity}.
\end{itemize}

Przepływ przetwarzania w tym wariancie przebiega następująco:
\begin{enumerate}
    \item klient wysyła żądanie HTTP do usługi \textit{Gateway},
    \item \textit{Gateway} przekazuje żądanie do usługi \textit{Search} przy użyciu gRPC,
    \item \textit{Search} wywołuje usługę \textit{Feature}, która oblicza wektor cech obrazu referencyjnego,
    \item \textit{Search} przekazuje otrzymany wektor do usługi \textit{Similarity},
    \item \textit{Similarity} porównuje wektor z indeksem i zwraca listę identyfikatorów obrazów,
    \item \textit{Search} przekazuje identyfikatory do \textit{Gateway},
    \item \textit{Gateway} wywołuje usługę \textit{Metadata}, która pobiera metadane dla wskazanych obrazów,
    \item \textit{Gateway} buduje końcową odpowiedź JSON i zwraca ją klientowi.
\end{enumerate}

Takie rozdzielenie odpowiedzialności pozwala lepiej odwzorować charakter systemów rozproszonych oraz umożliwia niezależne skalowanie wybranych komponentów. Jednocześnie wprowadza dodatkowy koszt związany z komunikacją pomiędzy usługami, serializacją wiadomości, zarządzaniem większą liczbą procesów oraz większą złożonością wdrożeniową.

Schemat infrastruktury wariantu mikroserwisowego przedstawiono na rysunku~\ref{fig:infra-microservices}, natomiast sekwencję przetwarzania w tym wariancie pokazano na rysunku~\ref{fig:seq-microservices}.

\begin{figure}[htbp]
    \centering
    \includegraphics[width=0.95\textwidth]{images/architecture/infra-microservices.pdf}
    \caption{Schemat infrastruktury wariantu mikroserwisowego}
    \label{fig:infra-microservices}
\end{figure}

\begin{figure}[htbp]
    \centering
    \includegraphics[width=0.95\textwidth]{images/architecture/microservices-sequence-diagram.pdf}
    \caption{Diagram sekwencji przetwarzania w wariancie mikroserwisowym}
    \label{fig:seq-microservices}
\end{figure}

\section{Stosy technologiczne}

Każdy z dwóch wariantów architektonicznych został zaimplementowany w trzech technologiach programistycznych. Takie podejście pozwala zmniejszyć ryzyko, że uzyskane wyniki będą specyficzne wyłącznie dla jednego środowiska uruchomieniowego.

\subsection{Java i Spring}

Implementacja w języku Java została oparta na frameworku Spring. Wariant monolityczny udostępnia REST-owy punkt końcowy wyszukiwania obrazów i realizuje wszystkie etapy przetwarzania w jednym procesie aplikacyjnym. Wariant mikroserwisowy wykorzystuje Spring jako bazę dla poszczególnych komponentów oraz biblioteki gRPC do komunikacji między usługami.

Java stanowi w badaniu reprezentanta środowiska o dojrzałym modelu wielowątkowości, rozbudowanym ekosystemie serwerowym oraz wysokiej przewidywalności zachowania w środowiskach produkcyjnych. Dzięki temu możliwe jest porównanie zachowania tej platformy z lżejszymi środowiskami interpretera i środowiska uruchomieniowego skryptowego.

\subsection{JavaScript i Node.js}

Wariant JavaScript został przygotowany w środowisku Node.js. Podobnie jak w przypadku Javy, przygotowano zarówno wersję monolityczną, jak i mikroserwisową z komunikacją gRPC. Zastosowanie Node.js pozwala uwzględnić w badaniu model wykonania oparty na pojedynczej pętli zdarzeń oraz nieblokującym przetwarzaniu operacji wejścia-wyjścia.

W kontekście badań porównawczych jest to istotne, ponieważ pozwala sprawdzić, jak architektura monolityczna i mikroserwisowa zachowują się w środowisku o odmiennym modelu współbieżności niż klasyczne środowisko JVM.

\subsection{Python i Flask}

Implementacja Python została oparta na frameworku Flask. Wariant monolityczny realizuje cały przepływ przetwarzania obrazu w jednej aplikacji, natomiast wariant mikroserwisowy rozdziela logikę pomiędzy pięć usług komunikujących się przez gRPC. Python włączono do badania jako reprezentanta środowiska szczególnie popularnego w zadaniach związanych z przetwarzaniem danych i obrazów.

Z punktu widzenia niniejszej pracy ważne było sprawdzenie, czy charakterystyka Pythona, w tym wyższy narzut interpretacyjny i odmienny model wykonania niż w Javie i Node.js, wpływa na skalę różnic pomiędzy monolitem a mikroserwisami.

\section{Dane testowe i przygotowanie środowiska}

We wszystkich implementacjach wykorzystano ten sam zbiór danych testowych. Obejmuje on:
\begin{itemize}
    \item obrazy źródłowe przechowywane w katalogu \texttt{data/images},
    \item wektory cech zapisane w pliku \texttt{data/features.json},
    \item metadane obrazów zapisane w pliku \texttt{data/metadata.json}.
\end{itemize}

Podczas uruchamiania aplikacji dane te są ładowane do pamięci operacyjnej. Takie podejście pozwala ograniczyć wpływ opóźnień pochodzących od zewnętrznych systemów bazodanowych oraz zapewnia spójne warunki wejściowe dla wszystkich testowanych wariantów. W konsekwencji pomiar koncentruje się na wpływie architektury i technologii implementacyjnej, a nie na narzucie związanym z warstwą persystencji.

\section{Znaczenie przyjętej implementacji dla badań eksperymentalnych}

Przyjęta konstrukcja aplikacji benchmarkowej ma istotne znaczenie metodologiczne. Po pierwsze, oba warianty architektoniczne realizują dokładnie ten sam proces biznesowy, co ogranicza wpływ czynników funkcjonalnych na wyniki eksperymentu. Po drugie, zastosowanie trzech języków programowania pozwala zweryfikować, czy zależności obserwowane pomiędzy monolitem a mikroserwisami są spójne w różnych środowiskach wykonawczych. Po trzecie, rozdzielenie wariantu mikroserwisowego na pięć komponentów powoduje, że w systemie pojawia się rzeczywisty koszt komunikacji rozproszonej, a więc cecha kluczowa z punktu widzenia celu pracy.

Przygotowana aplikacja benchmarkowa umożliwia zatem analizę trzech kluczowych obszarów:
\begin{itemize}
    \item wydajności, rozumianej jako czas odpowiedzi i przepustowość,
    \item skalowalności pionowej i poziomej,
    \item kosztów utrzymania wynikających z rozmiaru infrastruktury i liczby uruchamianych komponentów.
\end{itemize}

Tak zaprojektowane środowisko stanowi spójną i wiarygodną podstawę do dalszych badań eksperymentalnych przedstawionych w kolejnych rozdziałach pracy.

\section{Podsumowanie}

W rozdziale przedstawiono implementację aplikacji benchmarkowej wykorzystanej do porównania architektury monolitycznej i mikroserwisowej. Opisana aplikacja działa w domenie przetwarzania obrazów i realizuje proces wyszukiwania obrazów podobnych na podstawie ekstrakcji cech, porównania z indeksem oraz pobierania metadanych.

Aplikację zaimplementowano w trzech technologiach: Java z wykorzystaniem Spring, JavaScript w środowisku Node.js oraz Python z wykorzystaniem Flask. Dla każdej technologii przygotowano wersję monolityczną oraz wersję mikroserwisową opartą na pięciu usługach komunikujących się przez gRPC. Takie podejście zapewnia funkcjonalną równoważność badanych wariantów i tworzy właściwą podstawę do dalszej analizy wydajności, skalowalności i kosztów utrzymania.